%! TEX root = part5.tex
% the main TeX file which is intended to compile, :VimtexReload after adjustment
% :h vimtex-tex-root
\documentclass[../gbr_teknik.tex]{subfiles}
% \includeonlyframes{current}
% \setbeameroption{hide notes} % Only slides
% \setbeameroption{show only notes} % Only notes
% \setbeameroption{show notes on second screen=right} % Both

\begin{document}


\begin{frame}[mytitle=center,c,mybg=bglight1,light]

	\framesubtitle{Setelah mempelajari materi ini, mahasiswa mampu:}

	\begin{itemize}
		\item Menunjukkan macam-macam penyajian potongan
		\item Mendemonstrasikan cara menggambar potongan
		\item Mengklarifikasi benda yang tidak boleh dipotong
		\item Mendesain gambar benda dalam bentuk potongan
	\end{itemize}\end{frame}


\begin{frame}[mytitle=center,t,mybg=bglight2,light]\frametitle{Pengertian Gambar Potongan}

	\begin{itemize}
		\item Gambar potongan digunakan untuk memperjelas bagian dalam benda
		\item Menghindari penggunaan garis tersembunyi yang berlebihan
		\item Membantu pembaca memahami bentuk dan detail internal benda
	\end{itemize}\end{frame}


\begin{frame}[mytitle=standard,t,mybg=bglight3,light]\frametitle{Alasan Menggunakan Potongan}

	\begin{itemize}
		\item Garis tersembunyi terlalu banyak → membingungkan
		\item Bagian dalam sulit dipahami
		\item Potongan memberikan visual yang lebih jelas dan informatif
	\end{itemize}\end{frame}


\begin{frame}[mytitle=standard,c,mycolor=digiPH_uniblue,light]\frametitle{Macam-Macam Gambar Potongan}
	\begin{itemize}
		\item Potongan penuh (full section)
		\item Potongan setengah (half section)
		\item Potongan sebagian (partial section)
	\end{itemize}

	\scriptsize
	\begin{center}
		Catatan:

		Pemilihan jenis potongan disesuaikan dengan kebutuhan
		Untuk benda simetris → cukup potongan setengah
	\end{center}
	\normalsize\end{frame}

\begin{frame}[mytitle=standard,t,mycolor=digiPH_leaf,dark]\frametitle{Potongan Penuh}
	\begin{itemize}
		\item Seluruh benda dipotong
		\item Menampilkan seluruh bagian dalam
		\item Digunakan jika detail internal kompleks
	\end{itemize}\end{frame}

\begin{frame}[mytitle=standard,t,mycolor=digiPH_ocean,dark]
	\frametitle{Potongan Setengah}
	\begin{itemize}
		\item Digunakan untuk benda simetris
		\item Setengah tampak luar dan setengah tampak dalam
		\item Lebih efisien dalam waktu menggambar
	\end{itemize}
\end{frame}


\begin{frame}[mytitle=standard,t,mycolor=digiPH_gray,light]
	\frametitle{Potongan Sebagian}
	\begin{itemize}
		\item Hanya bagian tertentu yang dipotong
		\item Digunakan untuk detail lokal
		\item Lebih praktis dan tidak mubazir
	\end{itemize}
\end{frame}

\begin{frame}[mytitle=standard,t,mycolor=digiPH_ubcblue,dark]

	\frametitle{Cara Menggambar Potongan}
	\begin{itemize}
		\item Bagian yang terpotong diberi arsiran
		\item Sudut arsiran 45 derajat
		\item Jarak garis harus seragam
	\end{itemize}
\end{frame}

\begin{frame}[mytitle=standard,t,mycolor=digiPH_red,dark]
	\frametitle{Aturan Arsiran}
	\begin{itemize}
		\item Satu benda memiliki arah arsiran yang sama
		\item Benda berbeda memiliki arah arsiran berbeda
		\item Jarak arsiran disesuaikan ukuran gambar
	\end{itemize}
\end{frame}

\begin{frame}[mytitle=standard,t,mycolor=digiPH_cream,light]

	\frametitle{Benda Tipis}
	\begin{itemize}
		\item Tidak menggunakan arsiran biasa
		\item Digambar dengan garis tebal (di hitamkan)
		\item Antar bagian tetap diberi jarak
	\end{itemize}
\end{frame}

\begin{frame}[mytitle=standard,t,mycolor=digiPH_purple,dark]

	\frametitle{Benda yang Tidak Boleh Dipotong}
	\begin{itemize}
		\item Baut
		\item Paku keling
		\item Pasak
		\item Poros dan sirip
	\end{itemize}
\end{frame}

\begin{frame}[mytitle=standard,t,mycolor=digiPH_creamy,dark]

	\frametitle{Garis Pemotongan}
	\begin{itemize}
		\item Menggunakan garis strip titik
		\item Ujung garis dibuat tebal
		\item Dilengkapi anak panah
		\item Diberi huruf penunjuk
	\end{itemize}
\end{frame}

\begin{frame}[mytitle=standard,t,mycolor=digiPH_creamy,dark]
	\frametitle{Penunjukkan Potongan}
	\begin{itemize}
		\item Huruf pada garis potong harus sama dengan potongan
		\item Contoh: A-A, B-B
		\item Memudahkan identifikasi bagian
	\end{itemize}
\end{frame}

\begin{frame}[mytitle=standard,t,mycolor=digiPH_white,light]
	\frametitle{Kesimpulan}
	\begin{itemize}
		\item Gambar potongan memperjelas bagian dalam benda
		\item Terdapat beberapa jenis potongan
		\item Arsiran memiliki aturan khusus
		\item Tidak semua benda boleh dipotong
	\end{itemize}
\end{frame}

\begin{frame}[mytitle=standard,t,mycolor=digiPH_white,light]
	\includegraphics<1>[width=\textwidth]{../figures/part5_image1.jpeg}
	\includegraphics<2>[width=\textwidth]{../figures/part5_image2.jpeg}

	\includegraphics<3>[width=\textwidth]{../figures/part5_image3.jpeg}

	\includegraphics<4>[width=\textwidth]{../figures/part5_image4.jpeg}

	\includegraphics<5>[width=\textwidth]{../figures/part5_image5.jpeg}

	\includegraphics<6>[width=\textwidth]{../figures/part5_image6.jpeg}

	\includegraphics<7>[width=\textwidth]{../figures/part5_image7.jpeg}

	\includegraphics<8>[width=\textwidth]{../figures/part5_image8.jpeg}

	\includegraphics<9>[width=\textwidth]{../figures/part5_image9.jpeg}

	\includegraphics<10>[width=\textwidth]{../figures/part5_image10.jpeg}

	\includegraphics<11>[width=\textwidth]{../figures/part5_image11.jpeg}

	\includegraphics<12>[width=\textwidth]{../figures/part5_image12.jpeg}

	\includegraphics<13>[width=\textwidth]{../figures/part5_image13.jpeg}

	\includegraphics<14>[width=\textwidth]{../figures/part5_image14.jpeg}

	\includegraphics<15>[width=\textwidth]{../figures/part5_image15.jpeg}

	\includegraphics<16>[width=\textwidth]{../figures/part5_image16.jpeg}

	\includegraphics<17>[width=\textwidth]{../figures/part5_image17.jpeg}

	\includegraphics<18>[width=\textwidth]{../figures/part5_image18.jpeg}

	\includegraphics<19>[width=\textwidth]{../figures/part5_image19.jpeg}

	\includegraphics<20>[width=\textwidth]{../figures/part5_image20.jpeg}

	\includegraphics<21>[width=\textwidth]{../figures/part5_image21.jpeg}

	\includegraphics<22>[width=\textwidth]{../figures/part5_image22.jpeg}

	\includegraphics<23>[width=\textwidth]{../figures/part5_image23.jpeg}
	\includegraphics<24>[width=\textwidth]{../figures/part5_image24.jpeg}
\end{frame}


\end{document}
